\section{Stateful models}
\label{ch:stateful}

The sequence models keep a state per market, updated at every block from the stream of book
events, and ask the state for the prediction.  This section follows the stateful model of
\cref{sec:models:sequence} from its first version to its last, with, at each step, what changed, what it
gained and why.  Unless stated otherwise the numbers are those of the LOB-only model, with
no edges; the contribution of the edges is the subject of \cref{ch:graph}.

\subsection{First models}
\label{sec:stateful:first}

\Cref{tab:exp:retrain} gives the first numbers: the dynamic model, with or without the graph, beats
the blind baselines and momentum but loses to Mean Reversion by 3--8 points.  Adding
the lookback features of \cref{sec:models:lookback} to the window (without the direct
path) gains about one point at every horizon, and three seeds put the seed-to-seed
standard deviation at $0.001$--$0.002$.  The graph contributes nothing at 1h, $+0.1$ at
6h with mixed sign across seeds, and $+0.3$ at 12h with the same sign on all seeds.

\begin{table}[t]
\centering
\caption{Models before the direct path (seed 0; the lookback rows are the mean of three seeds).}
\label{tab:exp:retrain}
\begin{fittable}\begin{tabular}{lccc}
\toprule
Model & 1h & 6h & 12h \\
\midrule
Mean Reversion-$H$                    & 0.527 & 0.501 & 0.507 \\
CT-DGNN, graph                        & 0.496 & 0.471 & 0.430 \\
CT-DGNN, LOB-only                     & 0.492 & 0.466 & 0.430 \\
+ lookback 1/6/12h, graph (3 seeds)   & 0.5035 & 0.4769 & 0.4389 \\
+ lookback 1/6/12h, LOB-only (3 seeds) & 0.5037 & 0.4759 & 0.4358 \\
\bottomrule
\end{tabular}\end{fittable}
\end{table}

\subsection{Loss and imbalance variants}
\label{sec:stateful:variants}


With the LOB-only lookback model as reference, \cref{tab:exp:variants} compares the
fixed class weights with automatic (inverse-frequency) weights, with undersampling of
the \textsc{const} class during training, and with regression of the signed change
followed by a threshold calibrated on the last training day.  At 1h and 6h the fixed
weights are best; at 12h, where \textsc{const} is a third of the samples instead of a
half, undersampling gains five points.  Regression is far worse at every horizon: the
Huber loss is dominated by the few large moves and the threshold calibrated on
training data does not transfer.  The continuous-time anti-symmetric update
(CTAN, \cite{gravina2024ctan}) in place of the GRU gives no gain and is much worse
without the decay; the wallet behaviour features give $\pm0.3$ with no consistent sign.

\begin{table}[t]
\centering
\caption{Loss and imbalance variants, LOB-only model with lookback 1/6/12h, seed 0.}
\label{tab:exp:variants}
\begin{fittable}\begin{tabular}{lccc}
\toprule
Variant & 1h & 6h & 12h \\
\midrule
Focal, weights 2/1/2 (reference) & 0.504 & 0.476 & 0.437 \\
Focal, automatic weights          & 0.487 & 0.475 & 0.479 \\
Undersampling of \textsc{const}   & 0.486 & 0.467 & 0.488 \\
Regression + calibrated threshold & 0.410 & 0.428 & 0.410 \\
\bottomrule
\end{tabular}\end{fittable}
\end{table}

\subsection{Memory, direct path and lookbacks: the chain of improvements}
\label{sec:stateful:chain}


Starting from the lookback model, three additions were made one at a time and
evaluated at full length (\cref{tab:exp:chain,fig:stateful:chain}).  The hierarchical memory gains
$+1.2$ [$+0.9$, $+1.5$] at 1h and $+3.4$ [$+3.1$, $+3.7$] at 12h.  The direct path with
the global state gains $+2.9$ [$+2.6$, $+3.3$] at 1h and $+0.2$ [$0.0$, $+0.4$] at 12h.
Five lookbacks instead of three gain $+1.0$ [$+0.8$, $+1.1$] and $+1.9$ [$+1.7$,
$+2.0$].  The final LOB-only model is $+2.4$ [$+1.9$, $+2.8$] above Mean Reversion at 1h
and $-1.6$ [$-2.1$, $-1.1$] below it at 12h on all samples; on the deduplicated subset
it is above at both horizons ($+2.0$ and $+1.7$).  At 6h (\cref{tab:exp:6h}) it is
$+1.5$ [$+0.9$, $+2.1$] above on all samples and $+3.3$ above deduplicated.

\begin{table}[t]
\centering
\caption{Chain of improvements, LOB-only, seed 0, all samples.  Each $\Delta$ is the gain over the previous row with its 95\% hour-block bootstrap interval.}
\label{tab:exp:chain}
\begin{fittable}\begin{tabular}{lcccc}
\toprule
Configuration & 1h & $\Delta$ & 12h & $\Delta$ \\
\midrule
Mean Reversion-$H$                 & 0.5311 & & 0.5071 & \\
\midrule
lookback 1/6/12h                   & 0.5038 & & 0.4368 & \\
+ hierarchical memory (60 s, 1 h)  & 0.5161 & $+1.23$ [$+0.92$, $+1.50$] & 0.4706 & $+3.38$ [$+3.08$, $+3.67$] \\
+ direct path + global state       & 0.5452 & $+2.91$ [$+2.61$, $+3.25$] & 0.4727 & $+0.21$ [$+0.01$, $+0.39$] \\
+ five lookbacks (0.25/1/3/6/12 h) & 0.5547 & $+0.95$ [$+0.81$, $+1.08$] & 0.4914 & $+1.87$ [$+1.73$, $+2.01$] \\
\midrule
best $-$ Mean Reversion            & & $+2.36$ [$+1.93$, $+2.76$] & & $-1.57$ [$-2.13$, $-1.09$] \\
\bottomrule
\end{tabular}\end{fittable}
\end{table}

\begin{table}[t]
\centering
\caption{Best configuration at 6h (seed 0).}
\label{tab:exp:6h}
\begin{fittable}\begin{tabular}{lccc}
\toprule
Subset & Mean Reversion & LOB-only & graph \\
\midrule
all (9.3 M)           & 0.5069 & 0.5214 ($+1.45$ [$+0.89$, $+2.06$]) & 0.5279 ($+2.10$ [$+1.69$, $+2.51$]) \\
deduplicated (1.2 M)  & 0.5002 & 0.5333 ($+3.31$)                      & 0.5278 ($+2.76$) \\
extreme prices        & 0.4879 & 0.5136                                & 0.5154 \\
central prices        & 0.5199 & 0.5260                                & 0.5355 \\
\bottomrule
\end{tabular}\end{fittable}
\end{table}

\paragraph{Variants that did not help.}  Removing the decay ($0.475$ at 1h) or
replacing it with the fixed multi-scale decay ($0.509$, $0.427$) is worse than the
learned decay.  Initialising the direct path \emph{at} Mean Reversion instead of at
zero is $-2.1$ [$-2.7$, $-1.6$] at 1h; log-odds price features have no effect.  The
Mean-Reversion initialisation is also the clearest case in which a two-day short test
predicted the wrong sign ($+0.7$ and $+2.0$ in the short test, negative at full
length), which is why every claim in this section rests on full-length runs.

\paragraph{Seeds.}  Three seeds of the best LOB-only configuration give $0.5547$,
$0.5546$, $0.5549$ at 1h (deduplicated $0.538$--$0.544$): the gain over Mean Reversion
is stable to the third decimal.  At 12h they give $0.4914$, $0.4843$, $0.4695$
(standard deviation $0.009$), all below Mean Reversion on all samples and all above it
on the deduplicated subset: the 12h result is real on independent samples and fragile on
the full stream.

\begin{figure}[t]
\centering
\begin{tikzpicture}
\begin{axis}[thesis, width=0.47\textwidth, ybar, bar width=14pt, ymin=0.40, ymax=0.58, title={1 h}, title style={font=\footnotesize},
  ylabel={balanced accuracy, all samples}, symbolic x coords={lookback, + memory, + direct path, + 5 lookbacks}, xtick=data,
  x tick label style={rotate=30, anchor=east}, enlarge x limits=0.18, ytick={0.40,0.45,0.50,0.55}, name=left]
\addplot[fill=vizA, draw=white, line width=0.5pt, error bars/.cd, y dir=both, y explicit, error bar style={vizInk}]
  coordinates {(lookback,0.5038) (+ memory,0.5161) +- (0.0015,0.0015) (+ direct path,0.5452) +- (0.0017,0.0017) (+ 5 lookbacks,0.5547) +- (0.0007,0.0007)};
\addplot[vizB, dashed, line width=1pt, sharp plot] coordinates {(lookback,0.5311) (+ 5 lookbacks,0.5311)};
\end{axis}
\begin{axis}[thesis, width=0.47\textwidth, ybar, bar width=14pt, ymin=0.40, ymax=0.58, title={12 h}, title style={font=\footnotesize},
  symbolic x coords={lookback, + memory, + direct path, + 5 lookbacks}, xtick=data, x tick label style={rotate=30, anchor=east},
  enlarge x limits=0.18, ytick={0.40,0.45,0.50,0.55}, yticklabels={}, at={(left.east)}, anchor=west, xshift=1.2cm,
  legend style={at={(0.03,0.97)}, anchor=north west}]
\addplot[fill=vizA, draw=white, line width=0.5pt, error bars/.cd, y dir=both, y explicit, error bar style={vizInk}]
  coordinates {(lookback,0.4368) (+ memory,0.4706) +- (0.0015,0.0015) (+ direct path,0.4727) +- (0.001,0.001) (+ 5 lookbacks,0.4914) +- (0.0007,0.0007)};
\addplot[vizB, dashed, line width=1pt, sharp plot] coordinates {(lookback,0.5071) (+ 5 lookbacks,0.5071)};
\legend{stateful model, Mean Reversion}
\end{axis}
\end{tikzpicture}

\caption{The chain of improvements of the LOB-only stateful model at 1 h and 12 h (seed 0,
all samples; error bars are the 95\% hour-block bootstrap intervals of each step's gain).
The direct path is the large step at 1 h, the hierarchical memory at 12 h; at 12 h the
model stays below Mean Reversion on the full stream.}
\label{fig:stateful:chain}
\end{figure}

\subsubsection{Why the direct path matters}

Before the direct path, the model had to \emph{learn} mean reversion through a GRU over
window rows that contained the lookback change as one feature among fifty-five, and it
never fully did: it stayed three points below the rule.  With the direct path the rule is
a two-weight solution reachable by gradient descent, and the recurrent part only has to
learn the \emph{residual}: when not to revert.  The $+2.9$-point jump at 1h is the size
of what the recurrent part had been failing to express.  Initialising the path at the
rule instead of at zero is worse, which is instructive: starting from the rule puts the
model in a region where the gradient of the focal loss pushes the recurrent part toward
agreeing with the rule, and the residual is never learned; starting from zero lets both
parts grow together.

\subsubsection{Why the hierarchical memory matters, and why the decay stays short}

The learned decay rate converges to a 17-second half-life in every run.  This is not a
failure: at the two-second time scale of the stream, the per-block state is best used
as a sharp description of the last few updates, and a slow decay would blur it.  What
the fast state cannot do is remember the morning; the minute and hour states do, each
with its own GRU, and the gain is largest at 12h ($+3.4$ points), the horizon at which
the morning matters most.  A fixed multi-scale decay inside one state (the alternative
tested) is worse because it forces one set of weights to serve all time scales, which is
the problem the clockwork design avoids \cite{koutnik2014clockwork}.

\subsubsection{Why five lookbacks beat three}

The three lookbacks (1, 6, 12 h) describe the day; the two added ones (15 min and 3 h)
describe the recent past, where the dynamics of a prediction market are fastest.
A 15-minute change is close to the natural time scale of the one-tick oscillation and
lets the direct path revert \emph{that}, while reading the longer changes as trend.
The gain is consistent across seeds.

\subsection{Window encoders: space first, or time first}
\label{sec:stateful:encoders}


The order-book window of the best 1h model was encoded in three ways
(\cref{sec:models:encoders}): the reference GRU over the last 10 snapshots, the same
GRU over 32 snapshots (the longest that fits in memory) as a window-length control, and
the time-first encoder over 64 snapshots.  \Cref{tab:exp:encoders} compares the three
runs on the same validation samples.  The longer GRU window gains a quarter of a point;
the time-first encoder gains three quarters of a point more, to $0.5643$, the best
one-hour balanced accuracy of the study, one point above the reference with an interval
that excludes zero by a wide margin and thirty times the seed spread of the reference.
The gain is concentrated in the central price range ($+1.5$ points over the reference);
at extreme prices and on the deduplicated samples the three encoders are within their
intervals of each other, and the longer GRU is in fact the best on the deduplicated
samples.  The time-first encoder costs 14 milliseconds per block in the order-book
stage against 10 for the longer GRU and 5 for the reference.  A second seed and the
6h horizon (\cref{tab:exp:encoders:more}) confirm the result: seed 1 reproduces seed 0
to within a tenth of a point on every subset (two-seed means $0.5639$ against $0.5547$),
and at 6h the time-first encoder is $1.8$ points above the reference and $3.3$ above
Mean Reversion, this time on every subset, the deduplicated samples and the extreme
prices included, and above the best 6h graph run of \cref{sec:exp:graph}.

\begin{table}[t]
\centering
\caption{Three window encoders in the best 1h configuration, seed 0 (balanced accuracy;
differences in points with 95\% hour-block bootstrap intervals).}
\label{tab:exp:encoders}
\small
\begin{fittable}\begin{tabular}{lrcccccc}
\toprule
Subset & $n$ & MR & GRU $K{=}10$ & GRU $K{=}32$ & Time-first $K{=}64$ & $\Delta$ vs $K{=}10$ & 95\% CI \\
\midrule
All samples    & 9{,}864{,}125 & 0.5311 & 0.5547 & 0.5571 & 0.5643 & $+0.97$ & [$+0.63$, $+1.28$] \\
Deduplicated   & 1{,}255{,}651 & 0.5180 & 0.5382 & 0.5435 & 0.5410 & $+0.29$ & [$-0.06$, $+0.62$] \\
Extreme prices & 3{,}787{,}202 & 0.5230 & 0.5385 & 0.5420 & 0.5377 & $-0.07$ & [$-0.51$, $+0.34$] \\
Central prices & 6{,}076{,}923 & 0.5348 & 0.5602 & 0.5619 & 0.5748 & $+1.46$ & [$+1.15$, $+1.75$] \\
\bottomrule
\end{tabular}\end{fittable}
\end{table}

\begin{table}[t]
\centering
\caption{Time-first encoder against the reference GRU: seed 1 at 1h and seed 0 at 6h
(balanced accuracy; differences in points with 95\% hour-block bootstrap intervals).}
\label{tab:exp:encoders:more}
\small
\begin{fittable}\begin{tabular}{llrcccc}
\toprule
Horizon & Subset & $n$ & MR & GRU $K{=}10$ & Time-first $K{=}64$ & $\Delta$ [95\% CI] \\
\midrule
1h, seed 1 & All samples    & 9{,}864{,}125 & 0.5311 & 0.5546 & 0.5635 & $+0.89$ [$+0.54$, $+1.22$] \\
           & Deduplicated   & 1{,}255{,}651 & 0.5180 & 0.5437 & 0.5450 & $+0.13$ [$-0.19$, $+0.46$] \\
           & Extreme prices & 3{,}787{,}202 & 0.5230 & 0.5409 & 0.5403 & $-0.06$ [$-0.48$, $+0.34$] \\
           & Central prices & 6{,}076{,}923 & 0.5348 & 0.5588 & 0.5725 & $+1.37$ [$+1.04$, $+1.70$] \\
\midrule
6h, seed 0 & All samples    & 9{,}321{,}125 & 0.5069 & 0.5214 & \textbf{0.5396} & $+1.83$ [$+1.41$, $+2.21$] \\
           & Deduplicated   & 1{,}234{,}892 & 0.5002 & 0.5333 & 0.5370 & $+0.38$ [$+0.13$, $+0.64$] \\
           & Extreme prices & 3{,}744{,}385 & 0.4879 & 0.5136 & 0.5285 & $+1.49$ [$+0.94$, $+2.05$] \\
           & Central prices & 5{,}576{,}740 & 0.5199 & 0.5260 & 0.5468 & $+2.08$ [$+1.81$, $+2.36$] \\
\bottomrule
\end{tabular}\end{fittable}
\end{table}

Two caveats belong with this result.  The three momentum columns of the window are
computed over its $K$ snapshots, so they change meaning with $K$; and the snapshots are
on an event clock, one every 20 book updates, with no channel for the time elapsed
between them, so a 64-snapshot window covers seconds for a busy market and hours for a
quiet one.  The time gap between snapshots is the obvious channel to add next.

\subsubsection{Why reading time before space helps}
\label{sec:analysis:encoders}


The reference GRU condenses each book state into one vector before it sees any
sequence, so the temporal pattern of a single quantity (the ask depth draining over
the last minute, say) must survive a 55-to-128 compression at every step before the
recurrence can notice it.  The time-first encoder gives each quantity its own series
and a shared map from series to features, so such a pattern is represented directly,
and the attention across channels then decides which patterns matter together.  This
is the channel-independence argument of \textcite{nie2023patchtst} applied to an
order book, and the evidence that it is the order of operations and not the window
length is the control: a GRU over three times as many snapshots gets a quarter of the
gain.  The gain lives where the order-book series carry information, the central price
range, and vanishes at the extremes, where the book is one tick wide and the series
are noise; the deduplicated samples are dominated by the extremes
(\cref{sec:analysis:extreme}), which is why they show little.  The shared linear time
map also makes the encoder a cousin of the linear model of
\textcite{zeng2023transformers}: per channel it \emph{is} a linear model on the
normalised window, and the attention only mixes its outputs.

\subsection{The state given the control's inputs}
\label{sec:stateful:hand}
\label{sec:exp:handfeatures}

If the control of \cref{sec:stateless:control} wins because of what it reads and not because
of what it is, the stateful model should recover most of the gap once it reads the same
things.  The sixty-one features were therefore given to the best 1h model, the time-first
encoder with five lookbacks, as described in \cref{sec:models:handinputs}.
\Cref{tab:exp:handfeatures} gives the result on the model's own validation samples and
against the control on the control's samples.

\begin{table}[t]
\centering
\caption{The best 1h model with and without the control's features, seed 0 (balanced
accuracy; differences in points with 95\% hour-block bootstrap intervals).}
\label{tab:exp:handfeatures}
\small
\begin{fittable}\begin{tabular}{lcccccc}
\toprule
Subset & MR & Time-first & \textbf{+ features} & $\Delta$ vs time-first & LightGBM & $\Delta$ vs LightGBM \\
\midrule
All samples    & 0.5311 & 0.5643 & \textbf{0.5842} & $+1.99$ [$+1.65$, $+2.37$] & 0.6015 & $-1.73$ [$-2.28$, $-1.15$] \\
Deduplicated   & 0.5180 & 0.5410 & \textbf{0.5708} & $+2.97$ [$+2.78$, $+3.16$] & 0.5827 & $-1.20$ [$-1.57$, $-0.79$] \\
Extreme prices & 0.5230 & 0.5377 & \textbf{0.5667} & $+2.90$ [$+2.51$, $+3.29$] & 0.5804 & $-1.36$ [$-1.84$, $-0.79$] \\
Central prices & 0.5348 & 0.5748 & \textbf{0.5910} & $+1.62$ [$+1.28$, $+2.00$] & 0.6107 & $-1.97$ [$-2.58$, $-1.29$] \\
\bottomrule
\end{tabular}\end{fittable}
\end{table}

The features are worth two points on the full stream and three on the deduplicated and
extreme samples, the largest single gain of the study and the first one that is larger
on the independent samples than on the full stream: what the model lacked at the extreme
prices was the activity clock and the price level, and both are now in front of it.  The
gap to the trees is halved, not closed.  Calibration is not where the remaining
$1.7$ points are.  The model trained with the focal loss predicts a move for $53\%$ of
the samples against $45\%$ in the labels, but the post-hoc logit adjustment by the
training class prior \cite{menon2021logitadjustment}, which is consistent for the
balanced error when the probabilities are calibrated, is destructive here ($0.527$): the
focal weights already shift the probabilities towards the rare classes, and the
adjustment counts the shift twice.  The cross-fitted version, two logit offsets fitted on the
even validation hours and applied to the odd ones and vice versa, comes out at
$(-0.1, -0.1)$ and changes nothing ($0.5837$).  The decision rule of the model is already
balanced-accuracy-optimal to within a tenth of a point; what is missing is in the
function it has learned.

Two follow-up runs, each paired against the run with the features, exclude the two
obvious explanations of the residual.  Training with the control's objective, a
class-balanced cross-entropy in place of the focal loss, loses $0.47$ points
[$-0.74$, $-0.22$] on the full stream and $1.95$ [$-2.13$, $-1.78$] on the deduplicated
samples.  A predictor head four times wider and one layer deeper (256 units, two
layers), motivated by the 581 inputs that the default head maps through 64 units, loses
$0.37$ [$-0.70$, $-0.08$] and $1.90$ [$-2.12$, $-1.69$].  Both lose exactly where the
features had gained most, the deduplicated and the extreme samples, and tie on the
central prices; both predict fewer moves than the focal-loss run, a class mix within a
point of the labels', and on the sparse subsets that is what balanced accuracy
punishes.  With two logit offsets fitted crosswise on the validation hours the three runs
become indistinguishable ($0.5837$, $0.5848$, $0.5849$): the loss and the head moved the
decision bias of the model, not the function it learned.

The experiment of \cref{sec:exp:handfeatures} confirms the first reading and sharpens the
third.  Given the control's inputs, the recurrent model recovers half of the gap at once
(two points on the full stream, three on the independent samples), so most of what the
trees knew was indeed a matter of inputs.  The half that remains is not a matter of
calibration, since the decision rule is already optimal, and it is not obviously a matter
of capacity either: it is a difference between a model that sees every training row once,
in time order, through a 64-unit head, and six hundred rounds of boosting over the same
seven million rows.  Which of the two it is, the training regime or the head, is the
question the follow-up runs are meant to answer.

\subsection{The cost of the online regime}
\label{sec:stateful:regime}

A learner comparison on the control's own table, its 61 features, its training rows and
its validation samples (\cref{tab:exp:learners}), separates the candidate causes.  An MLP
of the shape of the wide head, trained as a tabular model is trained (shuffled
mini-batches, four passes, early stopping on 12 March), is $0.77$ points below the trees;
with a piecewise-linear encoding of each numeric feature over its training quantile bins
\cite{gorishniy2022embeddings} it is $0.41$ below.  The same MLP trained as the deep model
is trained, one pass over the rows in time order at the same learning rate, is $2.77$
below, and the deep model sits between the two.  Neither the learning rate nor the number
of steps explains it: a ten times larger rate in time order makes the one-pass MLP worse,
while one shuffled pass at the same rate recovers $0.8$ points and a second pass in time
order $0.6$.  The regime that the state imposes, one pass in time order, costs about two
points on this table; the learner itself costs less than one.  The combination of LightGBM
with the deep model's probabilities, a geometric mean with the weight chosen on the even
validation hours and applied to the odd ones and vice versa, is the first number of the
study above the trees, $+0.19$ [$+0.12$, $+0.27$] on all samples and $+0.23$ on the
deduplicated ones, whereas the same combination with the tabular MLP gives nothing
($+0.05$ [$-0.02$, $+0.11$]).  What the state adds is small, but it is there, and it is not
what another learner on the same table would add.

\begin{table}[t]
\centering
\caption{The control's table (61 features, training rows of 7--11 March, validation
samples of 13 March) under different learners, and the combination of LightGBM with other
models (geometric mean of the probabilities, weight cross-fitted on the validation hours).
Balanced accuracy; differences from LightGBM in points with 95\% hour-block bootstrap
intervals.}
\label{tab:exp:learners}
\small
\begin{fittable}\begin{tabular}{lcccc}
\toprule
Learner & All & $\Delta$ & Dedup. & $\Delta$ \\
\midrule
LightGBM & 0.6015 & & 0.5827 & \\
MLP, shuffled, four passes, early stopping & 0.5938 & $-0.77$ [$-0.99$, $-0.54$] & 0.5757 & $-0.70$ \\
MLP + piecewise-linear encoding & 0.5974 & $-0.41$ [$-0.65$, $-0.17$] & 0.5783 & $-0.44$ \\
MLP, one pass in time order (the deep model's regime) & 0.5738 & $-2.77$ [$-3.27$, $-2.18$] & 0.5623 & $-2.04$ \\
MLP, one pass in time order, rate $\times 10$ & 0.5675 & $-3.40$ [$-3.79$, $-2.97$] & 0.5355 & $-4.72$ \\
MLP, one pass shuffled & 0.5818 & $-1.97$ [$-2.34$, $-1.59$] & 0.5649 & $-1.79$ \\
MLP, two passes in time order & 0.5802 & $-2.13$ [$-2.62$, $-1.61$] & 0.5701 & $-1.26$ \\
Deep model + features & 0.5842 & $-1.73$ [$-2.27$, $-1.13$] & 0.5708 & $-1.20$ \\
\midrule
LightGBM $\oplus$ deep model + features & \textbf{0.6035} & $+0.19$ [$+0.12$, $+0.27$] & \textbf{0.5850} & $+0.23$ \\
LightGBM $\oplus$ MLP + piecewise-linear encoding & 0.6020 & $+0.05$ [$-0.02$, $+0.11$] & 0.5829 & $+0.02$ \\
\bottomrule
\end{tabular}\end{fittable}
\end{table}

\paragraph{Two passes and weight averaging.}
The two regime levers that the comparison suggests were then tried on the deep model
itself.  A second pass over the seven training days (states reset between the passes,
optimizer kept) gives $0.5907$, $1.08$ [$-1.56$, $-0.59$] below the trees on all samples
and $0.46$ [$-0.75$, $-0.14$] below on the deduplicated ones: the largest standalone
improvement after the features, and in the direction the tabular comparison predicted.
An exponential moving average of the weights along the pass (decay $0.998$ and $0.999$
per optimizer step, horizons of about $500$ and $1{,}000$ steps out of $3{,}000$), swapped
in before validation, does not help on its own: $0.5819$ and $0.5795$ on all samples, with
a decision bias toward DOWN that two cross-fitted offsets only partly remove ($0.5876$ and
$0.5859$), while on the deduplicated samples both tie with the trees ($0.5805$, $-0.22$
[$-0.68$, $+0.28$]).  What the averaged weights change is their complementarity: combined
with LightGBM as above they give $+0.68$ [$+0.57$, $+0.81$] and $+0.72$ [$+0.62$, $+0.84$]
on the deduplicated samples, three times the $+0.23$ of the run with features, and the
two-pass run combined gives $+0.26$ [$+0.16$, $+0.37$] on all samples.  The sequential
state is worth most exactly where the trees are weakest, on the sparse markets, and the
smoother weights carry that part better than the last ones.

\subsection{The state as the residual of the trees}
\label{sec:stateful:residual}

The residual model of \cref{sec:models:residual} stays paired with the run with features
(same loss, same seed, same samples).
The cross-fitted trees reproduce the control on its own samples once their fit rows follow the
control's rule (markets of the model's universe, last event at or after $t+H$): 0.6020 on all
samples and 0.5811 on the deduplicated ones against 0.6015 and 0.5827, with 94.6\% of the
decisions in common; fitted on every tensor row they were 1.6 points below, a reminder that the
population of the fit matters as much as the learner.

Table~\ref{tab:exp:residual} gives the result.  The residual model is the first run of the study
at or above the control on every subset: 0.6024 on all samples (+0.08, inside the interval) and
0.5862 on the deduplicated ones (+0.35 [+0.18, +0.51]); against the run with features the gain is
+1.5 to +2.0 points on every subset, paired, more than every other lever combined.  A second pass
over the training days moves the gain where the trees are weakest, +1.18 [+0.94, +1.42] on the
deduplicated samples and +0.50 [+0.21, +0.81] on the extreme ones, at no cost on all samples; with
two cross-fitted logit offsets the two-pass residual model is at or above the control on every
subset (0.6048 on all samples, +0.33 [$-0.01$, +0.64]).  This is
trees plus state, not a network beating the trees, and the number that measures the state is the
gain over the tree base (from 0.6020 / 0.5811 to 0.6024 / 0.5945), concentrated on the sparse and
extreme markets where a tabular model has least to work with.  The post-hoc combination of the
residual model with LightGBM adds a further +0.27 [+0.10, +0.43].

\begin{table}[t]
\centering
\caption{The deep model as the residual of the trees (1\,h, validation samples of the control).
Balanced accuracy; differences from LightGBM in points with 95\% hour-block bootstrap intervals;
``offsets'' = two logit offsets (DOWN, UP) fitted on the even validation hours for the odd ones and
vice versa.}
\label{tab:exp:residual}
\small
\begin{fittable}\begin{tabular}{lcccccc}
\toprule
Subset & LightGBM & Residual & $\Delta$ & Residual, 2 passes & $\Delta$ & 2 passes + offsets \\
\midrule
All      & 0.6015 & 0.6024 & $+0.08$ [$-0.17$, $+0.34$] & 0.6024 & $+0.09$ [$-0.24$, $+0.46$] & 0.6048 \\
Dedup.   & 0.5827 & 0.5862 & $+0.35$ [$+0.18$, $+0.51$] & \textbf{0.5945} & $+1.18$ [$+0.94$, $+1.42$] & 0.5933 \\
Extreme  & 0.5804 & 0.5817 & $+0.13$ [$-0.13$, $+0.40$] & \textbf{0.5854} & $+0.50$ [$+0.21$, $+0.81$] & 0.5865 \\
Central  & 0.6107 & 0.6112 & $+0.05$ [$-0.25$, $+0.37$] & 0.6093 & $-0.14$ [$-0.52$, $+0.33$] & 0.6123 \\
\bottomrule
\end{tabular}\end{fittable}
\end{table}

\subsubsection{The same calibration for everybody}
\label{sec:stateful:calibration}

The offsets are a post-hoc calibration that any model can have, and the comparison is only fair
when every model has it.  Applied to LightGBM they are worth +0.29 on all samples and +1.15 on the
deduplicated ones (0.6044 and 0.5943): most of what the residual runs gained on the sparse markets
was a better-placed decision threshold.  Calibrated against calibrated, the two-pass residual
model ties the trees on every subset (+0.04 [$-0.18$, +0.24] on all samples, $-0.10$ [$-0.26$,
+0.05] on the deduplicated ones), the one-pass residual model is at $-0.05$ and $-0.38$, and the
two-pass model without trees at $-1.42$.  At a one-hour horizon, then, the sequential state does
not add measurable balanced accuracy on top of calibrated trees; it adds a combination gain of the
same size as the calibration gain, which may be the same thing.  This is the result the thesis
reports at one hour, and the reason \cref{ch:external} looks for information the book does not
contain rather than for a better learner of the same table; \cref{ch:horizons} shows that the
answer is different at one minute.

\subsection{What the state is worth}
\label{sec:stateful:worth}

Read together, the section gives a precise account of the stateful family.  Its own
improvements are real and explicable: a direct path that contains the rule, a memory at three
rates, five lookbacks and an encoder that reads time before space take the LOB-only model
from eight points below Mean Reversion to three above it at one hour.  Its inputs were its
main handicap: the activity clock and the price level, which the event-clock tensor hides,
are worth two to three points once they are put in front of it.  Its training regime, one
pass in time order, costs about two points against the same learner trained as a tabular
model, and a second pass recovers a third of that.  And on top of calibrated trees that read
the same inputs, at one hour, the state adds nothing measurable on its own and a small
combination gain.  Where a state should matter is where a summary of the recent past is not
enough, the next minute of a quiet market, and \cref{ch:horizons} finds it exactly there.

